\documentclass[12pt,a4paper]{article}
\usepackage[utf8]{inputenc}
\usepackage[T1]{fontenc}
\usepackage[brazil]{babel}
\usepackage{geometry}
\usepackage{hyperref}
\usepackage{longtable}
\usepackage{array}

\geometry{margin=2.5cm}

\title{Booka: Sistema SaaS de Agendamento Multifuncional Inteligente}
\author{Centro Universitario Insted -- Analise e Desenvolvimento de Sistemas}
\date{2026}

\begin{document}

\maketitle

\section{Resumo}

O Booka e uma plataforma SaaS para organizacao de agendamentos entre clientes e prestadores de servicos. O sistema centraliza clientes, servicos, disponibilidade, bloqueios de agenda e agendamentos, reduzindo problemas de processos manuais como conflitos de horario, perda de informacoes e uso de planilhas.

\section{Objetivo}

O objetivo do projeto e propor uma solucao digital para apoio ao trabalho remoto ou hibrido, com foco em organizacao de processos, acompanhamento de tarefas, gestao de equipe e controle operacional.

\section{Metodologia}

O projeto utiliza Scrum, com backlog, sprints, planejamento, reunioes de acompanhamento, revisoes e retrospectivas. O GitHub Projects e usado como Kanban oficial.

\section{Equipe}

\begin{longtable}{|p{6cm}|p{8cm}|}
\hline
\textbf{Papel} & \textbf{Responsavel} \\
\hline
Product Owner / Scrum Master & Matheus Victor Moreira Yamanari \\
\hline
Frontend Angular/Ionic & Luis Fernando Franco \\
\hline
Frontend, telas e UI/UX & Thiago Almeida Akatsuka \\
\hline
Backend/API REST & Giulliano Ribeiro da Silva \\
\hline
Banco de dados, autenticacao e seguranca & Rubens Galina Junior \\
\hline
\end{longtable}

\section{Cronograma}

O planejamento considera seis meses: requisitos e arquitetura, banco/backend, frontend/integracao, funcionalidades centrais, testes/deploy e lockdown final.

\section{Conclusao}

O Booka consolida uma proposta academica e tecnica para digitalizacao de agendamentos, com organizacao de equipe por Scrum, controle por Kanban e documentacao formal em LaTeX.

\bibliographystyle{plain}
\bibliography{referencias}

\end{document}
